\subsection{Characterizations of the interval}

Lemma 6. --- Let D be a countable totally ordered set, not reduced to a single element, having a least and a greatest element. Suppose that D is without gaps (E, III, p.~73, exerc. 19), that is, every open interval {]}x, y{[}, where x and y are elements of D such that x \textless{} y, is nonempty. There then exists an isomorphism of ordered sets from I ∩ Q onto D.

Let \(a\) be the smallest element and \(b\) the largest element of D. By hypothesis, we have \(b \neq a\) and \(]a\) , \(x[\neq \emptyset\) for all \(x \in \mathrm{D} - \{a\}\) . The set \(\mathrm{D} - \{a\}\) is totally ordered, is nonempty, and has no least element; it is therefore infinite (E, III, p.~34, cor. 1 of prop. 3). The sets \(\mathbf{I} \cap \mathbf{Q}\) and D, infinite and countable, are equipotent to \(\mathbf{N}\) .

Let us choose bijections \(n \mapsto a_{n}\) and \(n \mapsto b_{n}\) from N onto \(I \cap Q\) and D respectively, such that \(a_{0} = 0\) , \(a_{1} = 1\) , \(b_{0} = a\) , \(b_{1} = b\). There exists a unique strictly increasing map \(f: I \cap Q \to D\) possessing the following properties: one has \(f(0) = a\) and \(f(1) = b\) ; for \(n \geqslant 2\) , we have \(f(a_{n}) = b_{m}\) , where m is the smallest natural integer for which the map from \(\{a_{0}, \ldots, a_{n}\}\) into D that coincides with f in \(\{a_{0}, \ldots, a_{n-1}\}\) and maps \(a_{n}\) to \(b_{m}\) is strictly increasing. These properties indeed define \(f(a_{n})\) by induction on n, the existence of the integer m being ensured by the fact that D has no gaps.

Since the map \(f\) is strictly increasing and \(\mathbf{I} \cap \mathbf{Q}\) is totally ordered, \(f\) defines an isomorphism of ordered sets from \(\mathbf{I} \cap \mathbf{Q}\) onto its image (E, III, p.~14, prop. 11). It remains for us to prove that \(f\) is surjective. To do this, let us prove by induction that \(b_{m}\) belongs to the image of \(f\) for all \(m \in \mathbf{N}\) .

We have \(b_{0}=f(0)\) and \(b_{1}=f(1)\) . Suppose that \(m\geqslant2\) and that, for \(0\leqslant k\leqslant m-1\) , there exists \(c_{k}\in I\cap Q\) such that \(f(c_{k})=b_{k}\) . We have \(c_{0}=0\) and \(c_{1}=1\) , since f is injective. Consider the smallest integer \(n\in N\) for which \(a_{n}\) does not belong to \(\{c_{0},\ldots,c_{m-1}\}\) and the map g from \(\{c_{0},\ldots,c_{m-1},a_{n}\}\) into D that coincides with f in \(\{c_{0},\ldots,c_{m-1}\}\) and maps \(a_{n}\) to \(b_{m}\) is strictly increasing; such an integer exists because \(I\cap Q\) is without gaps. Let \(f'\) be the map from \(\{a_{0},\ldots,a_{n}\}\) into D that coincides with f in \(\{a_{0},\ldots,a_{n-1}\}\) and maps \(a_{n}\) to \(b_{m}\). Let us prove that it is strictly increasing. Let \(j\in\{0,\ldots,n-1\}\) ; by definition of the integer n, there exists \(k\in\{0,\ldots,m-1\}\) such that \(a_{j}=c_{k}\) , or \(a_{j}<c_{k}\) and \(b_{m}\geqslant f(c_{k})\) , or \(a_{j}>c_{k}\) and \(b_{m}\leqslant f(c_{k})\) . Suppose \(b_{k}<b_{m}\) ; we then have \(g(c_{k})=f(c_{k})=b_{k}<b_{m}=g(a_{n})\) , whence \(c_{k}<a_{n}\) since g is strictly increasing, and then \(a_{j}\leqslant c_{k}<a_{n}\) ; moreover, \(f'(a_{j})=f(a_{j})\leqslant f(c_{k})=b_{k}<b_{m}=f'(a_{n})\) . Similarly, if \(b_{k}>b_{m}\) , it follows that \(a_{n}<c_{k}\leqslant a_{j}\) and \(f'(a_{j})=f(a_{j})\geqslant f(c_{k})=b_{k}>b_{m}=f'(a_{n})\) . Since f is itself strictly increasing, it follows that the map \(f'\) is strictly increasing.

If \(m'\) is the integer such that \(f(a_n) = b_{m'}\) , we have \(m' \leqslant m\) , by definition of f. If one had \(m' < m\) , one would have \(f(a_n) = b_{m'} = f(c_{m'})\) , whence \(a_n = c_{m'}\) , since f is injective, which contradicts the definition of \(a_n\) . Thus, \(m' = m\) and \(f(a_n) = b_m\) , which proves that \(b_m\) belongs to the image of f and completes by induction the proof of the surjectivity of f.

PROPOSITION 16. --- Let E be a totally ordered set not reduced to a single element. Suppose that every subset of E has a supremum and that there exists a countable subset of E that meets every open interval of the form {]}x, y{[}, where x and y are elements of E such that x \textless{} y. There then exists an isomorphism of ordered sets from I onto E.

Since ∅ and E each have a supremum in E, E has a least element a and a greatest element b. These are distinct since E is not reduced to a single element. Let D' be a countable subset of E that meets every open interval of E of the form {]}x, y{[} with x \textless{} y. Set D = D' ∪ \{a, b\}. The set D is totally ordered and has no gaps. By Lemma 6, there exists an order isomorphism f from I ∩ Q onto D.

For every \(t \in I\) , let \(g(t)\) be the upper bound of \(f([0, t] \cap \mathbf{Q})\) in E. For every \(x \in E\) , let \(h(x)\) be the upper bound of \(f^{-1}([a, x] \cap D)\) in I. The maps \(g: I \to E\) and \(h: E \to I\) thus defined are increasing, g coincides with f in \(I \cap Q\) and h coincides with \(f^{-1}\) in D.

We therefore have \(g(h(y)) = y\) for all \(y \in D\) . Let \(x \in E\) . If one had \(g(h(x)) > x\) , the interval {]}x, \(g(h(x))[\) would contain a point y of D and the relations \(g(h(y)) = y < g(h(x))\) would contradict the fact that \(g \circ h\) is increasing. Likewise, one does not have \(g(h(x)) < x\) . We therefore have \(g(h(x)) = x\) , which proves that \(g \circ h\) is the identity map of E. One shows similarly that \(h \circ g\) is the identity map of I. Thus, \(g: I \to E\) and \(h: E \to I\) are mutually inverse isomorphisms of ordered sets.

Remark. --- Let E be a totally ordered set. The set of open intervals of E (bounded or not) is stable under finite intersection. It is a base of a topology \(\mathcal{T}_{0}(\mathrm{E})\) on E (TG, I, p.~91, exerc. 5). The topology \(\mathcal{T}_{0}(\mathbf{I})\) is identical to the topology induced on \(\mathbf{I}\) by that of \(\mathbf{R}\) . It follows that, in prop. 16, every isomorphism of ordered sets from \(\mathbf{I}\) onto E is a homeomorphism of \(\mathbf{I}\) onto the topological space obtained by endowing E with the topology \(\mathcal{T}_{0}(\mathrm{E})\) .

COROLLARY. --- Let R be an equivalence relation in I. The following conditions are equivalent:

\begin{enumerate}
\def\labelenumi{(\roman{enumi})}
\item
  Every equivalence class modulo R is a closed interval of I, distinct from I;
\item
  There exists an increasing surjective map \(u\colon\mathbf{I}\to\mathbf{I}\) such that R is the equivalence relation associated with \(u\) .
\end{enumerate}

Such a map u, when it exists, is continuous and induces, by passage to the quotient, a homeomorphism of I/R onto I.

We shall denote by \(p: I \rightarrow I/R\) the canonical surjection.

Suppose condition (i) is satisfied. For A and B equivalence classes modulo R, write A \textless{} B if one has a \textless{} b for all \(a \in A\) and every \(b \in B\) . In I/R, the relation “A = B or A \textless{} B” is an order relation. Indeed, it is reflexive; it is antisymmetric because one cannot have simultaneously A \textless{} B and B \textless{} A; it is transitive because the relations A \textless{} B and B \textless{} C imply A \textless{} C. If A and B are distinct elements of I/R, they are disjoint closed intervals of I, and one then has either A \textless{} B or B \textless{} A. Endowed with the order relation thus defined, I/R is therefore totally ordered. The map \(p: I \to I/R\) is increasing.

The set I/R is not reduced to a single element, by virtue of (i).

Let F be a subset of I/R. Let us prove that F has a least upper bound in I/R. Put \(F' = \overline{p}^{-1}(F)\) ; denote by a the least upper bound of \(F'\) in I and by A the equivalence class of a under R. Since a majorizes \(F'\) in I, A majorizes F in I/R. Conversely, let \(B \in I/R\) be an upper bound of F; put \(b = \sup(B)\) . Every element of \(F'\) is then majorized by b. We therefore have \(a \leqslant b\) . Since the equivalence classes under R are closed intervals, a belongs to A and b to B. Consequently, we have \(A = p(a) \leqslant p(b) = B\) . This proves that A is the least upper bound of F.

Let A and B be elements of I/R such that A \textless{} B. Let a be the least upper bound of A and b the greatest lower bound of B. Since \(a \in A\) and \(b \in B\) , we have a \textless{} b. The equivalence class under R of any element of {]}a, b{[} is an element of the interval {]}A, B{[} of I/R. Since \(I \cap Q\) meets {]}a, b{[}, its image under p meets {]}A, B{[}.

We have thus proved that the totally ordered set I/R satisfies the hypotheses of prop. 16. There therefore exists an isomorphism of ordered sets \(f \colon \mathbf{I} / \mathrm{R} \to \mathbf{I}\) . The map \(u = f \circ p\) is a surjective and increasing map from I onto I and the equivalence relation associated with u is the relation R; this proves that condition (ii) is satisfied.

Suppose conversely that condition (ii) is satisfied. Let \(u: I \rightarrow I\) be an increasing and surjective map such that R is the equivalence relation associated with u.

Let \(a\) be a point of \(\mathbf{I}\) . The set \(\mathrm{A} = \overline{u}^{-1}(a)\) is an interval of \(\mathbf{I}\) , because the map \(u\) is increasing. Since \(u\) is surjective, A is neither empty nor equal to \(\mathbf{I}\) . Let \(b\) be the least upper bound of A in \(\mathbf{I}\) . We have \(u(b) \geqslant a\) . If \(u(b) > a\) , there exists \(c \in ]a, u(b)[\) and \(d \in \mathbf{I}\) such that \(u(d) = c\) , since \(u\) is surjective. Since \(u\) is increasing and \(a < u(d) < u(b)\) , \(d\) majorizes every element of A and we have \(d < b\) , which contradicts the hypothesis that \(b\) is the supremum of A. We therefore have \(u(b) = a\) , that is, \(b \in \mathrm{A}\) . Similarly, one shows that A contains its lower bound. The set A is therefore a closed interval of \(\mathbf{I}\) distinct from \(\mathbf{I}\) . This proves that condition (i) is satisfied.

We now prove that the map \(u\) is continuous. It suffices for this to show that, for every \(a \in \mathbf{I}\) , the sets \(\overline{u}^{-1}([a, \to[)\) and \(\overline{u}^{-1}([\leftarrow, a[)\) are open. Let \(b\) be the upper bound of \(\overline{u}^{-1}(a)\) ; we have \(u(b) = a\) . If \(x \in \mathbf{I}\) satisfies \(u(x) > a\) , one necessarily has \(x > b\) ; conversely, if \(x > b\) , we have \(u(x) \geqslant a\) and \(u(x) \neq a\) since \(b\) is the least upper bound of \(\overline{u}^{-1}(a)\) . We therefore have \(\overline{u}^{-1}([a, \to[) = ]b, \to[\) , which shows that the inverse image under \(u\) of the interval \(]a, \to[\) is open. One shows similarly that the inverse image of \(]\leftarrow, a[\) is open. It follows that the map \(u\) is continuous.

The map \(v \colon \mathbf{I} / \mathrm{R} \to \mathbf{I}\) deduced from \(u\) passing to the quotient is then continuous and bijective. Since \(\mathbf{I}\) is compact, \(\mathbf{I} / \mathrm{R}\) is quasi-compact (TG, I, p.~62, th. 2) and \(v\) is a homeomorphism (TG, I, p.~63, cor. 2).

PROPOSITION 17. --- Let X be a connected and compact topological space. Let a be a point of X, let U be a nonempty open and closed subset of \(X - \{a\}\).

\begin{enumerate}
\def\labelenumi{\alph{enumi})}
\item
  The closure \(\overline{U}\) of U in X is equal to \(U \cup \{a\}\) and is connected.
\item
  Let \(a'\) be a point of X distinct from \(a\) , let \(U'\) be a nonempty open and closed subset of \(X - \{a'\}\) . If \(a \notin U'\) and \(\overline{U} \cap \overline{U'} \neq \varnothing\) , we have \(\overline{U'} \subset U\) . Conversely, if \(a \in U'\) and \(X \neq U \cup U'\) , we have \(\overline{U} \subset U'\) .
\item
  There exists a point \(b\) of U such that \(X - \{b\}\) is connected.
\end{enumerate}

Let us prove assertion \(a\) ). Let V denote the complement of U in \(X - \{a\}\). Since U is closed in \(X - \{a\}\), V is open in \(X - \{a\}\) and a fortiori in X. We therefore have \(U \subset \overline{\mathrm{U}} \subset X - V = U \cup \{a\}\). Likewise, U is an open subset of X. Since X is connected, U is not closed in X, whence \(\overline{\mathrm{U}} = U \cup \{a\}\). Similarly, we have \(\overline{\mathrm{V}} = V \cup \{a\}\).

Let F and G be disjoint closed subsets of \(\overline{U}\) such that \(\overline{U} = F \cup G\) . Suppose that \(a \in G\) and let us prove that F is empty; we would reason similarly if \(a \in F\) . The set \(G \cup V = G \cup \overline{V}\) is a closed subset of X, disjoint from F, and we have \(X = \overline{U} \cup V = F \cup (G \cup V)\) . Since X is connected and G is nonempty, F is empty. This proves that \(\overline{U}\) is connected.

Let us prove \(b\) ). Suppose that \(a \notin \mathrm{U}'\) and that \(\overline{\mathrm{U}} \cap \overline{\mathrm{U}'} \neq \emptyset\) . By assertion \(a\) ), we have \(\overline{\mathrm{U}} = \mathrm{U} \cup \{a\}\) and \(\overline{\mathrm{U}'} = \mathrm{U}' \cup \{a'\}\) . Since \(a \notin \mathrm{U}'\) and \(a \neq a'\) , the sets U and \(\overline{\mathrm{U}'}\) have a common point. Still according to \(a\) ), \(\overline{\mathrm{U}'}\) is a connected subset of X; since it is contained in \(\mathrm{X} - \{a\}\) and meets the open and closed subset U, we have \(\overline{\mathrm{U}'} \subset \mathrm{U}\) , which is what had to be proved. The second assertion follows from the first by considering the complements in X of \(\overline{\mathrm{U}}\) and \(\overline{\mathrm{U}'}\) respectively.

Let us finally prove assertion \(c\) ). Suppose by contradiction that, for every \(x \in \mathrm{U}\) , the set \(\mathrm{X} - \{x\}\) is not connected, and choose subsets \(\mathrm{U}_x\) and \(\mathrm{V}_x\) , open and closed in \(\mathrm{X} - \{x\}\) , disjoint and nonempty, such that \(\mathrm{X} - \{x\} = \mathrm{U}_x \cup \mathrm{V}_x\) and \(a \in \mathrm{V}_x\) . By assertion \(b\) ), applied to the open and closed subsets \(\mathrm{U}\) , \(\mathrm{U}_x\) of \(\mathrm{X} - \{a\}\) and \(\mathrm{X} - \{x\}\) respectively, we have \(\overline{\mathrm{U}_x} \subset \mathrm{U}\) for all \(x \in \mathrm{U}\) . Let \(x\) and \(y\) be points of U such that \(x \in \mathrm{U}_y\) ; one therefore has \(x \neq y\) . Still by assertion \(b\) ), applied to the open and closed subsets \(\mathrm{U}_x\) and \(\mathrm{U}_y\) of \(\mathrm{X} - \{x\}\) and \(\mathrm{X} - \{y\}\) , the relation \(x \in \mathrm{U}_y\) implies the relation \(\overline{\mathrm{U}_x} \subset \overline{\mathrm{U}_y}\) . Consequently, the relations \(x \in \mathrm{U}_y\) and \(\overline{\mathrm{U}_x} \subset \overline{\mathrm{U}_y}\) are equivalent.

The set of subsets S of U such that \(\bigcap_{x\in\mathrm{S}}\overline{\mathrm{U}_{x}}\neq\varnothing\) is of finite character (E, III, p.~34), since the space X is compact. By E, III, p.~35, th. 1, there exists a maximal subset S of U such that \(\mathrm{C}=\bigcap_{x\in\mathrm{S}}\overline{\mathrm{U}_{x}}\) is not empty. Let \(c\) be a point of C. For every element \(x\) of S, we have \(c\in\overline{\mathrm{U}_{x}}\) , whence \(c\in\mathrm{U}\) ; then \(\overline{\mathrm{U}_{c}}\subset\overline{\mathrm{U}_{x}}\) . Consequently, one has \(\overline{\mathrm{U}_{c}}\subset\mathrm{C}\) ; then, by maximality of S, \(\mathrm{C}=\overline{\mathrm{U}_{c}}\) . For all \(x\in\mathrm{C}\) such that \(x\neq c\) , we also have \(\overline{\mathrm{U}_{x}}\subset\overline{\mathrm{U}_{c}}\) and \(\overline{\mathrm{U}_{x}}\neq\overline{\mathrm{U}_{c}}\) . By maximality of S, \(\mathrm{C}=\{c\}\) , therefore \(\overline{\mathrm{U}_{c}}=\{c\}\) , which contradicts the hypothesis that \(\mathrm{U}_{c}\) is a nonempty open and closed subset of \(\mathrm{X} - \{c\}\) .

COROLLARY. --- Let X be a compact connected topological space, and let N be the set of points of X whose complement is connected. The set X is the unique compact connected subset of X that contains N.

Let S be a connected and compact subset of X such that \(N \subset S\) . Suppose that \(S \neq X\) and let \(x \in X - S\) . By hypothesis, \(X - \{x\}\) is not connected; there therefore exist open and closed subsets U and V of \(X - \{x\}\) , disjoint and nonempty, such that \(X - \{x\} = U \cup V\) . We may assume that \(S \subset V\) . We have \(\overline{U} = U \cup \{x\}\) and \(\overline{V} = V \cup \{x\}\) , and these spaces are connected (III, p.~278, prop. 17, a)). By assertion c) of this proposition, there exists \(y \in U\) such that \(X - \{y\}\) is connected, which contradicts the inclusions \(N \subset S \subset V\) .

Lemma 7. --- Let T be a locally connected topological space. Let \(\mathcal{U}\) be a directed increasing family of open subsets of T, whose union is T. For \(t \in T\) and \(U \in \mathcal{U}\) , denote by \(U_{t}\) the connected component of t in U if \(t \in U\) and let us set \(U_{t} = \varnothing\) if \(t \notin U\) . For all \(t \in T\) , \(\bigcup_{U \in \mathcal{U}} U_{t}\) is the connected component of t in T.

For \(t \in \mathrm{T}\) , set \(\mathrm{C}_t = \bigcup_{\mathrm{U} \in \mathcal{U}} \mathrm{U}_t\) . We have \(t \in \mathrm{C}_t\) . The sets \(\mathrm{U}_t\) are open and connected, and the same holds for \(\mathrm{C}_t\) because we have \(t \in \mathrm{U}_t\) if \(\mathrm{U}_t \neq \varnothing\) (TG, I, p.~81, prop. 2). Let \(u\) and \(v\) be points of T such that \(\mathrm{C}_u \cap \mathrm{C}_v \neq \varnothing\) . Let \(t\) be a point of \(\mathrm{C}_u \cap \mathrm{C}_v\) ; let \(\mathrm{U} \in \mathcal{U}\) such that \(t \in \mathrm{U}_u\) and let \(\mathrm{V} \in \mathcal{U}\) such that \(v \in \mathrm{V}_v\) . Since \(\mathcal{U}\) is increasingly filtered, there exists \(\mathrm{W} \in \mathcal{U}\) which contains \(\mathrm{U} \cup \mathrm{V}\) . Then, \(\mathrm{W}_u \cap \mathrm{W}_v \neq \varnothing\) , therefore \(\mathrm{W}_u = \mathrm{W}_v\) . More generally, we have \(\mathrm{W}_u' = \mathrm{W}_v'\) for all \(\mathrm{W}' \in \mathcal{U}\) such that \(\mathrm{W} \subset \mathrm{W}'\) , therefore \(\mathrm{C}_u = \mathrm{C}_v\) . This shows that the sets of the form \(\mathrm{C}_t\) form a partition of T into connected open subsets. Consequently, for every \(t \in \mathrm{T}\) , \(\mathrm{C}_t\) is the connected component of \(t\) in T.

THEOREM 2. --- Let X be a compact connected topological space having a countable dense subset. Let a, b be distinct points of X. The following conditions are equivalent:

\begin{enumerate}
\def\labelenumi{(\roman{enumi})}
\item
  There exists a homeomorphism \(f\colon\mathbf{I}\to\mathbf{X}\) such that \(f(0)=a\) and \(f(1)=b\) ;
\item
  Every connected subset of X that contains \(\{a, b\}\) is equal to X;
\item
  The space X is locally connected, and every connected and compact subset of X that contains \(\{a,b\}\) is equal to X;
\item
  For every \(x \in X - \{a, b\}\) , the space \(X - \{x\}\) is not connected.
\end{enumerate}

Assertion (i) implies all the others.

Let \(x\) be a point of \(X - \{a, b\}\) . Since \(X - \{x\}\) contains \(\{a, b\}\) , assertion (ii) implies that this is not a connected subset of \(X\) , so that (ii) implies (iv).

Let us show that (iii) implies (iv). Suppose the hypotheses of (iii) are satisfied. Let \(x \in X - \{a, b\}\) and suppose that \(X - \{x\}\) is connected. Let T be the space \(X - \{x\}\) and let \(\mathcal{U}\) be the set of subsets of T of the form \(X - V\) , where V is a compact neighborhood of \(x\) . By Lemma 7, there exists a compact neighborhood V of \(x\) such that \(a\) and \(b\) belong to the same connected component of \(X - V\) . In particular, they belong to the same connected component of \(X - \mathring{V}\) ; this component is a compact, connected subset of X, distinct from X, which contradicts hypothesis (iii).

It remains to show that assertion (iv) implies condition (i).

Let \(x\) be a point of \(X - \{a, b\}\) and let U, V be open and closed, disjoint and nonempty subsets of \(X - \{x\}\) , such that \(X - \{x\} = U \cup V\) . By III, p.~278, prop. 17, c), there exists a point of U (resp. of V) whose complement in X is connected. Since X admits only two such points, \(a\) and \(b\) , one of them, say \(a\) , is contained in U and the other in V. By loc. cit., a nonempty open and closed subset \(U'\) of U contains a point of \(\{a, b\}\) , hence contains \(a\) . Applying this to \(U - U'\) , it follows that \(U = U'\) . Consequently, U is connected; it is the connected component of \(a\) in \(X - \{x\}\) . Likewise, V is the connected component of \(b\) in \(X - \{x\}\) .

For every \(x \in X - \{a, b\}\) , let \(U_x\) and \(V_x\) denote the connected components of a and b in \(X - \{x\}\) . By the preceding, they are open and closed in \(X - \{x\}\) , disjoint, and their union is equal to \(X - \{x\}\) .

Denote by \(\preccurlyeq\) the relation in X defined as follows: on the one hand, \(a \preccurlyeq x\) and \(x \preccurlyeq b\) for all \(x \in \mathrm{X}\) ; on the other hand, if \(x\) and \(y\) belong to \(\mathrm{X} - \{a, b\}\) , then \(x \preccurlyeq y\) if \(x \in \overline{\mathrm{U}_y}\) . For \(x\) and \(y\) in \(\mathrm{X} - \{a, b\}\) , the subsets \(\mathrm{V}_x\) and \(\mathrm{V}_y\) have the point b in common, and the relation \(x \preccurlyeq y\) is in fact equivalent to the assertion \(\overline{\mathrm{U}_x} \subset \overline{\mathrm{U}_y}\) (III, p.~278, prop. 17, b)). It follows that the relation \(\preccurlyeq\) is an order relation on X.

Let \(x\) and \(y\) be points of X such that one does not have \(x \preccurlyeq y\) . Necessarily, \(x \neq a\) and \(y \neq b\) , and \(x \in V_y\) . If \(x = b\) or \(y = a\) , we have \(y \preccurlyeq x\) . Suppose then that \(x\) and \(y\) are distinct from \(a\) and \(b\) . Since the parts \(U_x\) and \(U_y\) have the point a in common, we have the inclusion \(\overline{V_x} \subset \overline{V_y}\) (loc. cit.), whence, taking the closures of the complements, \(\overline{U_y} \subset \overline{U_x}\) and, a fortiori, \(y \preccurlyeq x\) . The relation \(\preccurlyeq\) in the space X is therefore a total order relation. For every \(x \in X - \{a, b\}\) , we moreover have \(U_x = ] \leftarrow, x[\) and \(V_x = ]x, \rightarrow[\) ; for \(x, y \in X - \{a, b\}\) , we have {]}x, y{[} = \(U_y \cap V_x\) . When x and y range over the points of \(X - \{a, b\}\) , the sets \(U_y \cap V_x\) , the sets \(U_y\) and the sets \(V_x\) form a base for a topology on X. Let \(\widetilde{X}\) denote the corresponding topological space and let i: \(X \to \widetilde{X}\) be the identity map of X. Since, for every \(x \in X\) , the sets \(U_x\) and \(V_x\) are open in X, the map i is continuous.

The space \(\widetilde{\mathbf{X}}\) is separated. Indeed, let \(x\) and \(y\) be distinct points of \(\widetilde{\mathbf{X}}\) such that \(x \preccurlyeq y\) . The sets \(]\leftarrow, y[\) and \(]x, \rightarrow[\) are open neighborhoods of \(x\) and \(y\) ; if they have a common point \(z\) , then \(]\leftarrow, z[\) and \(]z, \rightarrow[\) are disjoint open neighborhoods of \(x\) and \(y\) . Since X is compact, the map \(i\) is therefore a homeomorphism (TG, I, p.~63, cor. 2).

Therefore, the image under \(i\) of a countable dense subset of X intersects every nonempty open interval of the ordered set \((\mathrm{X}, \preccurlyeq)\) . It then follows from prop. 16 of III, p.~276 that there exists an isomorphism \(c\) of the ordered set \(\mathbf{I}\) onto \((\mathrm{X}, \preccurlyeq)\) . According to the remark following this proposition, this isomorphism is a homeomorphism of \(\mathbf{I}\) onto the topological space \(\widetilde{\mathrm{X}}\) . The map \(f = i^{-1} \circ c\) is then a homeomorphism of \(\mathbf{I}\) onto X which maps 0 to \(a\) and 1 to \(b\) .

